\chapter{Object Storage: MinIO and the S3 Model}
\label{ch:minio}

\section{The problem}

Assignment submissions are files --- PDFs, images, archives. Files in a
relational database bloat backups and stream badly; files on a pod's
local disk vanish with the pod. The industry's answer is \emph{object
storage}: an HTTP service storing immutable blobs in named
\emph{buckets} under string \emph{keys}, metadata included, no
directories (the \code{/} in a key is naming convention, not structure).
Amazon S3 defined the API; MinIO is a self-hostable server speaking the
same API --- so course-service uses the standard AWS SDK and cannot tell
the difference.

\section{The configuration, annotated}

From \code{configurations/course-service.yml}:

\begin{lstlisting}[language=yaml]
app:
  s3:
    endpoint: ${APP_S3_ENDPOINT:http://localhost:9000}         (*@\co{1}@*)
    public-endpoint: ${APP_S3_PUBLIC_ENDPOINT:http://localhost:9000} (*@\co{2}@*)
    access-key: ${APP_S3_ACCESS_KEY:minioadmin}                (*@\co{3}@*)
    secret-key: ${APP_S3_SECRET_KEY:minioadmin}
    bucket: ${APP_S3_BUCKET:ts-submissions}
    presigned-url-expiry-minutes: ${APP_S3_PRESIGNED_EXPIRY_MIN:15} (*@\co{4}@*)
\end{lstlisting}

\begin{description}
  \item[\co{1}] Where \emph{the service} talks to MinIO. In the cluster:
  \code{http://minio:9000}.
  \item[\co{2}] Where \emph{the browser} talks to MinIO --- the subject of
  the pitfall below.
  \item[\co{3}] Credentials with dev defaults; the cluster injects real
  ones from the \code{minio-app} Secret. \code{app.s3.*} is custom
  (\code{@ConfigurationProperties}) config, not Spring-managed --- the
  same override machinery applies regardless.
  \item[\co{4}] How long a presigned link stays valid.
\end{description}

\section{Presigned URLs: the pattern worth knowing}

Naive file download routes bytes \emph{through} the service:
browser $\rightarrow$ service $\rightarrow$ MinIO and back, tying up an
application thread per transfer. The S3-native alternative: the service
uses its credentials to \emph{sign a URL} --- bucket, key, expiry, HMAC
signature --- and hands the URL to the browser, which downloads
\emph{directly} from storage. The service spends microseconds; storage
does the streaming; the link self-expires. The same trick works for
uploads.

\begin{pitfall}[the two-endpoints problem]
Symptom: uploads work, but clicking a download link fails with
\code{ERR\_NAME\_NOT\_RESOLVED} on \code{minio:9000}. Cause: a presigned
URL embeds the endpoint \emph{it was signed against}. The service signs
against \code{minio:9000} --- a name only resolvable inside the cluster
network --- and the browser lives outside it. Hence two configured
endpoints: sign download links against \code{public-endpoint}, talk to
storage over \code{endpoint}. Every containerized S3 setup meets this
bug once; some meet it again with the signature itself, because the
signature covers the host --- you cannot sign against one host and
rewrite the URL to another afterwards.
\end{pitfall}

\begin{hardware}
MinIO here is a single container with a single PVC --- fine for a
homelab's submissions, nothing like production MinIO (multi-node erasure
coding) or S3's eleven nines of durability. The interesting part is that
course-service cannot tell: the same SDK calls would hit real S3 with
only the endpoint and credentials changed. That is the point of speaking
a standard API.
\end{hardware}

\section{Outside the project}

The same service on AWS proper: set \code{endpoint} and
\code{public-endpoint} to the real S3 URL (or leave the SDK's default),
replace static keys with an IAM role assumed by the pod (IRSA on EKS ---
no long-lived credentials at all), and turn on bucket versioning and
lifecycle rules (auto-expire old submissions to cheaper storage). The
concepts --- buckets, keys, presigning, the identity that signs ---
transfer unchanged; what changes is who manages durability and how
identity is proven.

\begin{quiz}
\begin{enumerate}
  \item Why do presigned URLs scale better than proxying downloads
  through the service, and what new failure mode do they introduce?
  \item Explain why \code{endpoint} and \code{public-endpoint} must
  differ inside the cluster, and why post-signing URL rewriting is not a
  fix.
  \item A presigned URL leaks in a chat log. What limits the damage, by
  design?
  \item What changes --- and what does not --- when this code moves from
  MinIO to real S3?
\end{enumerate}
\end{quiz}
